\scanpage{137}
% archive foliation: 136
expliciter cette action en tant qu'action sur $\mathcal D$ (défini par (22)).

Si une structure complexe sur $\mathbb R^2$ est définie par l'élément $z\in\mathcal D$, $z=\unsure{x}+iy$, alors l'application \inserted{$\unsure{\mathbb R|\mathbb C}$-linéaire} $\varphi:L_z=\mathbb C(z,1)\longrightarrow V$ s'identifie à:
\[
\underset{\substack{\Vert\\a+ib}}{\lambda}\longmapsto
\operatorname{Re}\bigl(\deleted{\unsure{x}}\underbrace{\lambda(ze_1+e_2)}_{[(ax-by)+i(ay+bx)]e_1+(a+ib)e_2}\bigr)
=(ax-by)e_1+ae_2
\]
les images inverses $\lambda_1,\lambda_2$ de $e_1,e_2$ dans celle-ci sont données par
\[
\lambda_1=-\frac{i}{y}\qquad
\lambda_2=1+i\frac{x}{y}
\]
d'où
\[
\lambda_2/\lambda_1=i(y+ix)=-x+iy
\]
donc on a
\[
\underbrace{e_2/e_1}_{\substack{\text{quotient dans }\mathbb C^*\\
\text{au sens de la}\\
\text{structure compl. de }\mathbb R^2\\
\text{définie par }z=x+iy\in\mathcal D}}
=\overline{(-x+iy)}=-x-iy=-z
\quad\in\mathcal D^-.
\]

Donc si une structure de Torelli sur $E$ est définie par la \deleted{\unsure{donnée}}\inserted{base} $(z_1,z_2)$ de $\Pi\subset t_E$, alors l'élément $z\in\mathcal D$ correspondant est donné par
\begin{equation}
\deleted{z_2/z_1=-x+iy=-\overline z\quad\text{où }z=x+iy.}
\qquad\inserted{z=-(z_2/z_1)}
\tag{24}
\end{equation}

Soit donc
\[
g\in\mathrm{Sl}(2,\mathbb Z)\qquad
g=\begin{pmatrix}a&b\\c&d\end{pmatrix}
\qquad
\begin{gathered}
a,b,c,d\in\mathbb Z\\
ad-bc=1
\end{gathered}
\]
et soit \deleted{\unsure{un u}} $u:\mathbb Z^2\xrightarrow{\sim}\Pi$ une structure de Torelli, correspondant à
